% s01-intro.tex —— 导引（原文 p.1–3，无编号节）
\section*{导引}
\addcontentsline{toc}{section}{导引}
\markright{导引}
\srcpage{1}

我们三人中的一位，Langlands，受我们中的第二位 Ngô 关于基本引理的工作的鼓舞——二十多年间基本引理一直没有证明，这在许多方面阻碍了自守形式解析理论的任何实质性进展——曾勾画过一个用以建立函子性的纲领，而函子性是这一理论的两个主要目标之一。第三位，Frenkel，注意到，从对应的几何形式（有时称为互反律或 Langlands 对应）——即自守形式与 Galois 表示之间预期存在的对应——中提取的一些想法与公式，有力地支持了这一策略。这些就是本文的三个出发点。

基本引理现已借助 [N] 的结果得到确立\yizhu{Ngô Bao Châu 于 2008 年宣布证明基本引理（\emph{Le lemme fondamental pour les algèbres de Lie}），本文文献表 [N] 即该文，Publ. Math. IHÉS 111 (2010)。}{依据：文献 [N] 条目；「2008 年宣布」为学界公认事实（B 级），非原书 p.1 内容}，因此稳定迹公式已近在咫尺。
尽管我们清楚仍然存在的障碍及其困难，也还没有能力克服它们，我们在本文及后续文章中着手 [L4] 所勾画纲领的第一步，其目标是建立自守形式与代数几何（丢番图的或非丢番图的）之间预期的联系。我们清楚，第一步——更不用说其后各步——在大多数（几乎是全部）算术学家乃至自守形式理论专家看来并不现实。但我们仍希望他们以开放的心态仔细阅读本文及后续文章。

尽管在这第一篇文章中我们只考虑数域以及有限域上一维函数域上的理论，我们希望在后续工作中，不是以统一的方式，而是以平行的方式，处理复代数曲线上的理论；不过总还有一些模糊之处有待澄清。\footnote{Edward Frenkel 的研究由 DARPA（资助号 HR0011-09-1-0015）与巴黎数学科学基金会资助。Ngô Bao Châu 的研究由 Simonyi 基金会资助。}

\srcpage{2}
我们可以用几句话说明本文的内容。在第一节中，我们简要概述 [L3] 的想法，但除自守 L 函数的对数导数之外，还使用这些函数本身。我们注意到，这种处理函子性的方式带有经典类域论的余味。第二节是唯一将域 $F$ 限制为函数域的一节；在这一节我们验证一个几何引理，其合理性来自对如下对应的预见：$\mathbb{F}_q$ 上某条曲线的函数域 $F$ 的自守形式，与 $\operatorname{Gal}(F^{\mathrm{sep}}/F)$ 的 $\ell$-进表示之间的对应。接着，在第三节回顾并重新表述稳定迹公式之后，我们在第四节引入所谓的迹公式的阿代尔化。迹公式中会出现测度。过去这些测度使公式的有效使用变得困难，因为它们引入一些因子，妨碍了 Poisson 方法的应用，而后者几乎是唯一看来有希望的方法。然而有一个重要的观察：在完全无害的假设下，公式稳定化形式中最有趣的部分，即对正则元素求和的部分，差不多就是对一个 $F$ 上有限维向量空间 $V$ 求和，被求和的是定义在 $V\otimes\mathbb{A}_F$ 上的一个函数；该函数的行为使得可以对 $V\subset V\otimes\mathbb{A}_F$ 这一对使用 Poisson 公式。不过我们必须使用 Kottwitz 关于测度所给因子的先前工作，因为他的公式让我们能断定这些因子是常数——我们认为这是一个奇迹。请读者留意我们对 Kottwitz 工作的引用以及我们由之得出的结论，并加以思考。我们强调，这些因子只在稳定公式中才会简化。我们还要补充，为了使用 Poisson 公式，我们将需要截断空间 $V\otimes\mathbb{A}_F$；为此我们借用了 Jayce Getz 的一个想法，他在另一情境中使用这一想法，但始终在迹公式的框架内。我们感谢他为我们解释这一想法。

这种处理迹公式所提出的解析问题的可能性，其重要性无法过分强调。迄今为止还没有人注意到这一点。不过本文不讨论 Poisson 和的分析，而是在最后一节验证：其中出现的变换 $\hat\theta$ 在 $F$ 上离散和的第一项 $\hat\theta(0)$ 是主贡献，即一维自守表示的贡献。这是一个很好的征兆。但我们还要补充，这一贡献虽然占主导，却远不是最重要的。

\srcpage{3}
我们中的一位特别高兴本文出现在献给 Paulo Ribenboim 的一期中，因为正是应他的请求，内窥的最初几次尝试才被写下来并发表在《加拿大数学杂志》上。
